\chapter{Pengukuran dan Keliling}\label{ch:g4:measure}

Kelas 3 mengukur dengan penggaris, timbangan dan jam
(\cref{ch:g3:measure}); bab ini menambahkan tabel satuan --- mesin
yang mengubah segalanya --- serta dua cara mengukur sebuah bangun
datar: panjang tepinya (\omterm{def:g4:measure:perimeter}{keliling}) dan ruang di dalamnya (\omterm{def:g4:measure:area}{luas}, yang
untuk sementara dibilang dalam satuan petak).

\section{Tabel satuan}

\begin{method}[Mengubah satuan dengan tabel]\label{met:g4:measure:table}
Untuk panjang, setiap kolom bernilai sepuluh kali kolom berikutnya:
\begin{center}
\renewcommand{\arraystretch}{1.3}
\begin{tabular}{c|c|c|c|c|c|c}
km & hm & dam & m & dm & cm & mm \\
\hline
 & & & $3$ & $4$ & $0$ & \\
\end{tabular}
\end{center}
\begin{enumerate}
  \item Tulislah bilangannya satu angka per kolom, dengan angka
        satuan ukurannya di kolom satuannya (di sini $3.4$ m: angka
        $3$ di bawah m);
  \item bacalah bilangan itu dalam satuan yang baru: $3.4$ m
        $= 340$ cm (isilah kolom yang kosong dengan \omterm{def:g1:counting:zero}{nol});
  \item \omterm{def:g2:measure:units}{massa} (kg\,\dots\,g) dan \omterm{def:g2:measure:units}{kapasitas} (L\,\dots\,cL) bekerja
        dengan cara yang sama.
\end{enumerate}
\end{method}

\begin{example}\label{ex:g4:measure:convert}
$5$ km $= 5\,000$ m; \quad $270$ cm $= 2.7$ m (angka $2$ mendarat di
bawah m); \quad $3$ kg $250$ g $= 3\,250$ g; \quad
$2.5$ L $= 250$ cL.
\end{example}

\section{Keliling}

\begin{definition}[Keliling]\label{def:g4:measure:perimeter}
\emph{Keliling}\index{keliling} sebuah \omterm{def:g4:geometry:polygon}{segi banyak} adalah panjang
seluruh tepinya: jumlahkan panjang semua sisinya (dalam satuan yang
sama!). Jalan pintas untuk bangun yang klasik:
\[
\text{persegi panjang: } P = 2 \times (L + w),
\qquad
\text{persegi: } P = 4 \times c .
\]
\end{definition}

\begin{example}\label{ex:g4:measure:perimeter}
Sebuah kebun berbentuk persegi panjang $12$ m kali $7$ m:
\[
P = 2 \times (12 + 7) = 2 \times 19 = 38 \text{ m}
\]
pagar. Sebuah bingkai foto persegi dengan sisi $18$ cm:
$P = 4 \times 18 = 72$ cm.
\end{example}

\section{Luas: membilang petak}

\begin{definition}[Luas dengan membilang]\label{def:g4:measure:area}
\emph{Luas}\index{luas} sebuah bangun yang digambar di kertas
berpetak adalah banyaknya petak yang ditutupinya. Dua setengah petak
dihitung sebagai satu. Luas menjawab ``berapa besar permukaannya'',
\omterm{def:g4:measure:perimeter}{keliling} menjawab ``berapa panjang tepinya'' --- dua pertanyaan yang
berbeda!
\end{definition}

\begin{omfigure}
\begin{tikzpicture}[scale=0.55]
  \draw[gray!40, thin] (0,0) grid (13,4);
  \draw[very thick, omDef, fill=omDef!15] (1,1) rectangle (5,3);
  \node[font=\small, omDef] at (3,2) {$8$ petak};
  \draw[very thick, omThm, fill=omThm!15]
    (7,1) -- (11,1) -- (11,2) -- (9,2) -- (9,3) -- (7,3) -- cycle;
  \node[font=\small, omThm] at (8.4,1.5) {$6$ petak};
\end{tikzpicture}

{\small Dua bangun di kertas berpetak: luasnya $8$ dan $6$ petak.
Sekarang telusuri tepinya lalu bilanglah satuannya: $12$ untuk
keduanya! \omterm{def:g4:measure:perimeter}{Kelilingnya} sama, luasnya berbeda.}
\end{omfigure}

\begin{example}[Setengah petak]\label{ex:g4:measure:half}
Sebuah segitiga siku-siku yang digambar di kertas berpetak dengan
sisi siku-siku $4$ dan $2$ menutupi $3$ petak penuh dan $2$ setengah
petak\,\dots\ bilanglah dengan teliti: luasnya adalah setengah dari
persegi panjang $4 \times 2$, yaitu $4$ petak. Memotong persegi
panjang sepanjang diagonalnya selalu memberi dua segitiga yang sama
luasnya.
\end{example}

\begin{example}[Luas sama, keliling berbeda]\label{ex:g4:measure:samearea}
Persegi $4 \times 4$ dan persegi panjang $8 \times 2$ sama-sama
menutupi $16$ petak --- luasnya sama. Tepinya: $16$ satuan untuk
persegi, $20$ untuk persegi panjang. Tidak ada ukuran yang
menentukan ukuran yang lain: bangun yang panjang dan tipis punya
tepi yang banyak untuk permukaan yang sedikit.
\end{example}

\section{Lama waktu}

\begin{method}[Menjumlahkan waktu yang melewati jam bulat]\label{met:g4:measure:time}
Untuk menjumlahkan atau mengurangkan lama waktu, kerjakan dengan
lompatan melewati jam yang bulat (ingat $1$ jam $= 60$ menit):
\[
9\text{ jam }40 + 35 \text{ menit}:
\quad 9\text{ jam }40 \xrightarrow{+20} 10\text{ jam }00
\xrightarrow{+15} 10\text{ jam }15 .
\]
Jangan pernah menjumlahkan menit melewati $60$ tanpa mengubahnya:
$40 + 35 = 75$ menit $= 1$ jam $15$ menit.
\end{method}

\begin{example}\label{ex:g4:measure:timetable}
Keretanya berangkat pukul $8{:}47$ dan tiba pukul $11{:}05$.
Lamanya: $8{:}47 \to 9{:}00$ adalah $13$ menit;
$9{:}00 \to 11{:}00$ adalah $2$ jam; $11{:}00 \to 11{:}05$ adalah
$5$ menit. Seluruhnya: $2$ jam $18$ menit.
\end{example}

\section{Latihan}

\begin{exercise}[$\star$]\label{exo:g4:measure:1}
Ubahlah dengan tabel satuan: $7$ m menjadi cm; $3.2$ km menjadi m;
$85$ mm menjadi cm; $4\,500$ m menjadi km.
\end{exercise}

\begin{exercise}[$\star$]\label{exo:g4:measure:2}
Ubahlah: $2$ kg menjadi g; $1\,800$ g menjadi kg dan g; $3.5$ L
menjadi cL; $25$ cL menjadi L.
\end{exercise}

\begin{exercise}[$\star$]\label{exo:g4:measure:3}
Hitunglah \omterm{def:g4:measure:perimeter}{kelilingnya}: persegi panjang $9$ cm $\times$ $5$ cm;
persegi dengan sisi $7.5$ cm; segitiga dengan sisi $6$ cm, $8$ cm dan $10$ cm.
\end{exercise}

\begin{exercise}[$\star$]\label{exo:g4:measure:4}
Sebuah persegi panjang punya \omterm{def:g4:measure:perimeter}{keliling} $26$ cm dan panjang $8$ cm.
Carilah lebarnya (carilah dulu panjang $+$ lebar).
\end{exercise}

\begin{exercise}[$\star$]\label{exo:g4:measure:5}
Di kertas berpetak, gambarlah bangun berbentuk huruf L yang menutupi
tepat $10$ petak, lalu bilanglah satuan \omterm{def:g4:measure:perimeter}{kelilingnya}.
\end{exercise}

\begin{exercise}[$\star$]\label{exo:g4:measure:6}
Gambarlah dua persegi panjang berbeda yang luasnya $12$ petak.
Hitunglah \omterm{def:g4:measure:perimeter}{keliling} masing-masing. Mana yang lebih ``bulat'', mana yang lebih ``tipis''?
\end{exercise}

\begin{exercise}[$\star$]\label{exo:g4:measure:7}
Sebuah segitiga siku-siku digambar di kertas berpetak dengan sisi
siku-siku $6$ dan $3$. Berapa luasnya dalam satuan petak? (Pakailah \cref{ex:g4:measure:half}.)
\end{exercise}

\begin{exercise}[$\star$]\label{exo:g4:measure:8}
Hitunglah lama waktunya: dari $10{:}20$ sampai $12{:}45$; dari
$7{:}35$ sampai $9{:}10$. Sebuah film berdurasi $1$ jam $50$ dan
mulai pukul $20{:}30$: pukul berapa ia berakhir?
\end{exercise}

\begin{exercise}[$\star$]\label{exo:g4:measure:9}
Kirana berlari $3$ putaran mengelilingi lapangan berbentuk persegi
panjang $60$ m kali $45$ m. Berapa meter yang ia tempuh?
\end{exercise}

\begin{exercise}[$\star\star$]\label{exo:g4:measure:10}
Seorang petani ingin memagari sawah berbentuk persegi dengan sisi
$85$ m, dan menyisakan gerbang $4$ m tanpa pagar. Berapa meter pagar yang harus ia beli?
\end{exercise}

\begin{exercise}[$\star\star$]\label{exo:g4:measure:11}
Benar atau salah, dengan contoh di kertas berpetak: ``jika dua
bangun punya \omterm{def:g4:measure:perimeter}{keliling} yang sama, luasnya sama''; ``jika dua bangun
punya \omterm{def:g4:measure:area}{luas} yang sama, \omterm{def:g4:measure:perimeter}{kelilingnya} sama''.
(\Cref{ex:g4:measure:samearea} dan gambar pada bab ini akan membantu.)

\end{exercise}
